\subsection*{Masking for Multi-Levels}

Thread masking can be applied at different parallel levels. The following example demonstrates this:

\begin{lstlisting}[language=croqtile,caption={Masking for Multi-Levels}]
__co__  void foo() {
  parallel p by 3, q by 4 {
    inthreads (p < 2 && q == 0) {
      // divergent code for thread (0, 0), (1, 0)
      // thread-mask: 011-0001
    }
    inthreads (q == 1) {
      // divergent code for thread (0, 1), (1, 1), (2, 1)
      // thread-mask: 111-0010
    }
  }
}
\end{lstlisting}

This code uses two levels of parallelization:

\begin{itemize}[leftmargin=1.6em]
\item Outer level (\texttt{p}): Parallel count of 3
\item Inner level (\texttt{q}): Parallel count of 4
\end{itemize}

The first \texttt{inthreads} predicate (\texttt{p \textless{} 2 \&\& q == 0}) restricts execution to:

\begin{itemize}[leftmargin=1.6em]
\item Outer parallel IDs: \texttt{0, 1} (excluding 2)
\item Inner parallel ID: \texttt{0} only
\end{itemize}

The thread mask (\texttt{011-0001}) reflects these constraints, where the two levels are separated by a hyphen (-).

The second predicate (\texttt{q == 1}) imposes no restriction on the outer level (\texttt{p}), allowing all outer IDs (\texttt{0, 1, 2}) while restricting the inner level to \texttt{q == 1}.
